\chapter{Introduction}
\label{ch:introduction}

Large Language Models (LLMs) have achieved remarkable fluency across diverse tasks, demonstrating an unprecedented ability to synthesize information and generate human-like text. However, their pervasive vulnerability to \textbf{hallucinations}—generating fluent, authoritative statements that contradict real-world factual ground truth—remains the primary barrier to their reliable deployment in high-stakes domains such as medicine, law, finance, and scientific discovery.

% ******************************************************************************
\section{Problem Statement}
\label{sec:problem_statement}
% ******************************************************************************

Existing approaches to hallucination detection predominantly treat verification as an undifferentiated neural sequence matching task. Standard Retrieval-Augmented Generation (RAG) and LLM-as-a-Judge frameworks retrieve unstructured web or Wikipedia passages and task cross-encoders with evaluating semantic entailment.

This monolithic approach suffers from three fundamental structural vulnerabilities:
\begin{enumerate}
    \item \textbf{Attention Dilution \& Lexical Overlap Bias:} When generated prose shares extensive vocabulary with a retrieved passage but alters a single critical parameter (e.g., asserting an Apple M3 Max has 40 million transistors instead of 92 billion), continuous embedding attention distributes across the matching nouns, causing neural judges to misclassify the contradiction as supported.
    \item \textbf{Absence of Arithmetic and Symbolic Calculators:} Neural language models are fundamentally probabilistic token predictors rather than symbolic arithmetic engines; they cannot subtract values, check numerical error margins, or verify interval containment with mathematical guarantees~\cite{fisch2026finground}.
    \item \textbf{Prohibitive Latency and Compute Footprint:} Multi-agent search evaluators such as Google SAFE~\cite{wei2024safe} require repeated web queries and multi-billion parameter LLM calls, incurring over 10 seconds of latency and substantial financial costs per response.
\end{enumerate}

% ******************************************************************************
\section{The HalDetect Hypothesis \& Research Questions}
\label{sec:research_questions}
% ******************************************************************************

We hypothesize that \textbf{fact verification is fundamentally a triage problem rather than a monolithic generation task}. By parsing unstructured text into typed propositions and routing them to specialized deterministic Arithmetic Logic Units (ALUs) and structured knowledge graphs, fact-checking systems can achieve near-perfect factual precision, eliminate lexical overlap false alarms, and operate in real time on standard consumer hardware.

To systematically investigate this hypothesis, this thesis addresses four core Research Questions:
\begin{itemize}
    \item \textbf{RQ1 (Routing Efficacy):} Does routing propositions into specialized modality pipelines outperform monolithic neural cross-encoders on balanced factual benchmarks?
    \item \textbf{RQ2 (Factual Precision):} Can domain-specific neuro-symbolic ALUs eliminate false alarms in quantitative calculations, temporal tenure, and academic citation attribution?
    \item \textbf{RQ3 (Generative Generalization):} How effectively does HalDetect catch real-world factual confabulations generated by frontier foundation models in open-ended prose?
    \item \textbf{RQ4 (Computational Efficiency):} What are the latency, memory, and energy trade-offs of neuro-symbolic triage compared to neural RAG and cloud LLM judges?
\end{itemize}

% ******************************************************************************
\section{Summary of Contributions}
\label{sec:contributions}
% ******************************************************************************

The primary contributions of this thesis are as follows:
\begin{enumerate}
    \item \textbf{The HalDetect Neuro-Symbolic Verification Architecture:} We design and implement a hierarchical verification pipeline that combines a sub-$0.05\text{ms}$ deterministic taxonomy router with five domain-specific ALUs: a Dimensional SI Quantitative ALU, a 7-Continent Topological Superlative Resolver, an Entity-Attribute Statement Matcher, a CrossRef Citation Gateway, and a Calibrated Relational NLI fallback.
    \item \textbf{Paragraph-Level Anaphora Normalization \& Conjunction DAG:} We introduce a forward antecedent coreference tracker that binds orphaned pronoun chains ($< 0.1\text{ms}$) and a multi-aspect DAG that enforces 3-valued conjunction logic under open-world assumptions.
    \item \textbf{The Symmetrical Master Benchmark Corpus ($N=1,000$):} We construct a perfectly balanced, 500-TRUE / 500-FALSE benchmark across five modalities, establishing an evaluation standard that renders trivial class-prior exploitation ineffective.
    \item \textbf{Comprehensive Multi-LLM Road Test \& Empirical Validation:} We evaluate 1,798 propositions generated by seven foundation models (Qwen 2.5, Mistral, DeepSeek-R1, Gemma 31B, GPT-OSS 120B, MiniMax M3), proving HalDetect achieves \textbf{100.0\% recall on real hallucinations} (catching 31/31 errors) where Modern SOTA RAG and MiniCheck-7B achieve 0.0\% recall ($p < 10^{-60}$).
\end{enumerate}

% ******************************************************************************
\section{Thesis Organization}
\label{sec:thesis_organization}
% ******************************************************************************

The remainder of this dissertation is organized as follows: Chapter~\ref{ch:related_work} reviews related literature across atomic fact-checking, neural NLI, RAG, and attention probing. Chapter~\ref{ch:methodology} formalizes the HalDetect architecture, its 5-tuple schema atomizer, and specialized ALUs. Chapter~\ref{ch:experiment} details dataset provenance, baseline configurations, and statistical protocols. Chapter~\ref{ch:results} presents comprehensive empirical results across laboratory benchmarks, generative road tests, and ablations. Chapter~\ref{ch:discussion} synthesizes findings, provides theoretical failure analyses, and details threats to validity. Chapter~\ref{ch:conclusion} concludes with future directions and released open-source artifacts.
